\documentclass[conference]{IEEEtran}

\usepackage{cite}
\usepackage{amsmath,amssymb,amsfonts}
\usepackage{graphicx}
\usepackage[hyphens]{url}
\Urlmuskip=0mu plus 1mu\relax
\usepackage{textcomp}
\usepackage{xcolor}
\usepackage{booktabs}
\usepackage{tagging}

\usepackage{sc26repro}

% The next two lines MUST stay commented for final submission.
%\usetag{explanation}
%\usetag{example}

\begin{document}

\twocolumn[%
{\begin{center}
\Huge
Appendix: Artifact Description
\end{center}}
]

%%%%%%%%%%%%%%%%%%%%%%%%%%%%%%%%%%%%%%%%%%%%%%%%%%%%%
%  AD Appendix
%%%%%%%%%%%%%%%%%%%%%%%%%%%%%%%%%%%%%%%%%%%%%%%%%%%%%

\appendixAD

\section{Overview of Contributions and Artifacts}

\subsection{Paper's Main Contributions}

\begin{description}
\item[$C_1$] We provide the first large-scale, trace-grounded
characterization of AI network sensitivity across realistic
workloads, fabrics, and scales up to 4{,}096 GPUs. We publicly
release our trace dataset (\textasciitilde2~TB of trace data on up to
4k GPUs), separate the latency, bandwidth and jitter contributions,
and identify the regimes where each dominates the bottleneck.
\item[$C_2$] We develop and validate a compositional, linear-programming--based
trace analysis methodology with signature-based caching that makes
production-scale sensitivity studies feasible while preserving the
dependency structure, regime boundaries, and sensitivity
classifications of the underlying workload.
\item[$C_3$] We derive closed-form latency--bandwidth performance
envelopes that expose safe relaxation regions and danger zones, and
translate workload behavior into actionable provisioning guidance
for AI-supercomputer fabrics.
\end{description}

\subsection{Computational Artifacts}

This paper is accompanied by two computational artifacts. Artifact
$A_1$ contains the source code of the Composite-LP toolchain together
with the packaged sensitivity-sweep outputs and the figure-generation
scripts needed to regenerate every figure in the paper from the
packaged data. Artifact $A_2$ is the companion collection of
execution traces (GOAL format) used throughout the study.

\begin{description}
\item[$A_1$] \texttt{https://github.com/tommasobo/\allowbreak SC\_Tracing.git}
\item[$A_2$] \texttt{http://storage2.spcl.ethz.ch/\allowbreak traces/ai/}
\end{description}

Digital Object Identifiers for both artifacts will be minted at
the AE stage, as allowed by the SC26 AD/AE policy.

\begin{center}
\small
\begin{tabular}{rll}
\toprule
Artifact & Contributions &  Related \\
ID       & Supported     &  Paper Elements \\
\midrule
$A_1$   &  $C_1$, $C_2$, $C_3$ & Figs.\ 1, 3--10 \\
\midrule
$A_2$   &  $C_1$               & Table~II, Figs.\ 1, 3--10 \\
\bottomrule
\end{tabular}
\end{center}


%%%%%%%%%%%%%%%%%%%%%%%%%%%%%%%%%%%%%%%%%%%%%%%%%%%%%%%%
\section{Artifact Identification}
%%%%%%%%%%%%%%%%%%%%%%%%%%%%%%%%%%%%%%%%%%%%%%%%%%%%%%%%

\newartifact

\artrel

Artifact $A_1$ ships every component needed to regenerate the
paper's figures from raw GOAL traces, organized as the following
subtrees:

\begin{itemize}
\item \texttt{solver/} --- the per-collective LP solver and
  program-level composition code.
\item \texttt{tools/LogGOPSim/} --- the LogGOPSim source with a
  Makefile and a build-script wrapper.
\item \texttt{pipeline/} --- thin drivers (\texttt{run\_composite\_lp.py},
  \texttt{run\_lgs.py}, \texttt{demo.py}) that compose the solver
  and LGS into one pipeline.
\item \texttt{data/} --- the precomputed sensitivity CSVs, so
  figure reproduction is decoupled from solver-time requirements,
  plus a tiny demo GOAL trace for the pipeline demo.
\item \texttt{scripts/} --- one figure-generation script per
  paper figure; invoked by \texttt{reproduce\_all.py}.
\end{itemize}

The toolchain ingests GOAL dependency graphs produced from NCCL
traces (contribution $C_1$), decomposes the per-iteration workload
into unique collective signatures, and builds a per-signature
parametric linear program that it solves in a single dual-simplex
sweep over latency or bandwidth; a lightweight program-level LP
then composes the per-signature piecewise-affine slopes into an
end-to-end sensitivity curve (contribution $C_2$). The sweep
outputs are the envelopes and regime boundaries shown in Figures
1, 3--10 (contribution $C_3$). Figure 2 is a schematic of this
pipeline and has no reproducible output.

\artexp

The results should demonstrate three main claims of the paper.
First, our Composite-LP methodology reproduces the end-to-end
latency and bandwidth sensitivity of the Monolithic-LP baseline
within approximately one percent on every workload for which
Monolithic is feasible, while using orders of magnitude less memory
and solve time. This memory advantage is what makes the
sensitivity analysis feasible at production scales: the Monolithic
LP would require tens of TB of RAM at 4{,}096 GPUs, whereas the
Composite LP runs in roughly 200\,MB.

Second, the derived latency-bandwidth envelopes should reveal
qualitatively distinct regimes across workloads. Llama 70B training
is expected to be bandwidth-dominated below ${\sim}200$\,Gbps with a
latency cliff only at tens of microseconds; Grok 314B at
4{,}096 GPUs is expected to be compute-bound across most of the
studied range, with latency mattering only well above the regime of
modern datacenter fabrics; vLLM Llama 8B inference is expected to
be memory-bound with very little sensitivity to the network at all.

Third, the jitter analysis on Grok 314B is expected to show that
the workload absorbs sparse jitter with sub-percent slowdown but
degrades sharply once a large fraction of messages is perturbed,
supporting the provisioning conclusions drawn in the paper.

All figures reported in the paper were produced by the same scripts
included in $A_1$.

\arttime

We distinguish three levels of reproduction.

\emph{Figure reproduction} (from packaged CSVs; the default path
for reviewers).

\begin{itemize}
\item Artifact setup: \textasciitilde2 min to clone the repository
  and install Python dependencies via
  \texttt{pip install -r requirements.txt}.
\item Artifact execution: \textasciitilde1 min to run
  \texttt{python3 reproduce\_all.py} on a standard laptop; each
  individual figure script completes in under 15 s.
\item Artifact analysis: immediate. Figures are written to
  \texttt{figures/} as PDF and PNG; visual comparison with the paper
  is sufficient.
\end{itemize}

\emph{Pipeline demo} (builds LogGOPSim from source and replays the
shipped demo GOAL trace through the LGS stage).

\begin{itemize}
\item Additional setup: \textasciitilde30 s on first invocation to
  compile LogGOPSim via \texttt{pipeline/build\_tools.sh}.
\item Additional execution: \textasciitilde10 s to run
  \texttt{python3 reproduce\_all.py --pipeline} end-to-end (the
  LGS replay of the demo GOAL is dominated by txt2bin conversion
  and runs in well under a second per latency point).
\end{itemize}

\emph{Data reproduction (from raw GOAL traces in $A_2$ to the
packaged sensitivity CSVs in $A_1$).} Running the end-to-end
Composite-LP pipeline from traces to CSVs requires a workstation
with a licensed Gurobi installation. We report approximate
end-to-end wall times below (on an AMD EPYC host with 128 cores and
380\,GiB of RAM):

\begin{itemize}
\item Trace download from $A_2$: dominated by the user's bandwidth
  (\textasciitilde 2\,TB total across all workloads; individual
  workloads are 1--10\,GiB for Llama / vLLM and up to
  \textasciitilde 1.3\,TiB for Grok 314B at 4{,}096 GPUs).
\item GOAL parsing and signature extraction: \textasciitilde 5 min for
  Llama~70B, \textasciitilde 2 min for vLLM, \textasciitilde 30 min
  for Grok 314B.
\item Composite-LP sweep over latency or bandwidth:
  \textasciitilde 1--5 min per workload; Grok 4{,}096 GPUs finishes
  in under 10 min on a single socket.
\item Monolithic-LP cross-check (only feasible for small scales):
  \textasciitilde 83 min at Llama 8B / 16 GPUs; infeasible beyond
  \textasciitilde 1{,}024 GPUs due to memory (see Figure~7 of the
  paper).
\end{itemize}

Reviewers who opt for the default figure-reproduction path do not
need to execute any part of the data-reproduction stage.

\artin

\artinpart{Hardware}

Figure reproduction requires no specialized hardware: any x86 or
ARM machine with ${\sim}1$\,GiB of free RAM is sufficient. The
upstream composition and solver stages used to produce the packaged
CSVs used an AMD EPYC host (128 cores, 380\,GiB RAM) with Gurobi;
scale-out tracing of Grok 314B (up to 4{,}096 GPUs) was performed on
the Microsoft Azure ND GB200 v6 cluster, and Llama / vLLM tracing
was performed on the CSCS Alps supercomputer (NVIDIA GH200,
Slingshot-11). None of that infrastructure is needed to reproduce
the figures from the packaged data.

\artinpart{Software}

Figure reproduction requires Python 3.8 or later and three
widely-available Python packages:

\begin{itemize}
\item \texttt{matplotlib} \textgreater= 3.5
\item \texttt{numpy} \textgreater= 1.21
\item \texttt{pandas} \textgreater= 1.3
\end{itemize}

A \texttt{requirements.txt} is shipped with the artifact.

The optional \texttt{--pipeline} path additionally builds
LogGOPSim from its bundled source, which requires
\texttt{g++}, \texttt{gengetopt}, and \texttt{re2c} (Ubuntu:
\texttt{apt-get install g++ gengetopt re2c}).

Regenerating the Composite-LP or Monolithic-LP sensitivity CSVs
from raw GOAL traces additionally requires a Gurobi installation
(version 10.0 or later). Gurobi is commercial software, but the
free academic license offered at \url{gurobi.com/academia} is
sufficient to reproduce every LP solve reported in the paper.
The Gurobi dependency applies \emph{only} to the LP stage of
data reproduction; the figure-reproduction path and the LGS
pipeline demo do not invoke any solver.

\artinpart{Datasets / Inputs}

All inputs required to regenerate the figures are shipped inside
$A_1$ under \texttt{data/}: ${\sim}6$\,MiB of CSV files containing the
latency and bandwidth sensitivity sweeps, memory footprints, and
hardware reference measurements for the three workloads.
The full collection of raw execution traces is published separately
as $A_2$ (\textasciitilde 2\,TB, GOAL format). $A_2$ is not required
for figure reproduction; it is the input to the LP solver stage that
generated the packaged CSVs in $A_1$.

\artinpart{Installation and Deployment}

Clone the repository (the artifact lives on the \texttt{main}
branch, which is the default), create a Python virtual environment,
and install the three runtime dependencies:

\begingroup\small
\begin{verbatim}
git clone -b main \
  https://github.com/tommasobo/SC_Tracing.git
cd SC_Tracing
python3 -m venv .venv
source .venv/bin/activate
pip install -r requirements.txt
\end{verbatim}
\endgroup

No compilation is required. The figure-reproduction path uses
only Python and the packaged CSVs under \texttt{data/}.

\artcomp

The end-to-end path from cloning the repository to producing
every paper figure is three commands:

\begingroup\small
\begin{verbatim}
git clone -b main \
  https://github.com/tommasobo/SC_Tracing.git
cd SC_Tracing
pip install -r requirements.txt
python3 reproduce_all.py
\end{verbatim}
\endgroup

Concretely, the workflow has three stages:

\begin{itemize}
\item $T_1$ --- Install Python dependencies with
  \texttt{pip install -r requirements.txt}.
\item $T_2$ --- Run \texttt{python3 reproduce\_all.py}, which
  iterates over the paper's figure list and invokes one Python
  script per figure. Each script reads CSV files from
  \texttt{data/} and writes the corresponding PDF and PNG to
  \texttt{figures/}.
\item $T_3$ (optional) --- Run
  \texttt{reproduce\_all.py --pipeline}, which first builds
  LogGOPSim from source under \texttt{tools/LogGOPSim/} and
  replays the shipped demo GOAL trace through the LGS stage,
  then proceeds with $T_2$. This exercises the same LGS
  machinery used at paper scale. Expected wall time: ${<}1$\,min
  on top of the figure-reproduction path.
\end{itemize}

Task dependencies: $T_1 \rightarrow T_2$ and $T_1 \rightarrow T_3
\rightarrow T_2$. No parameter sweeps or randomness are involved
in $T_2$: every figure is deterministic given the packaged CSVs,
so repetition is not required. Users may reproduce a subset of
figures by passing their paper numbers via \texttt{--only}, e.g.\
\texttt{python3 reproduce\_all.py --only 3 4 5}.

\emph{Driving the pipeline on production-scale traces.} To run the
Composite-LP sweep or LGS replay on the actual paper workloads,
download one of the per-workload trees from $A_2$ (each ships a
GOAL file plus a \texttt{comm\_dep.csv} that disambiguates
sends/recvs) and invoke:

\begingroup\small
\begin{verbatim}
python3 pipeline/run_composite_lp.py \
  --goal <workload>/output.goal \
  --comm-dep <workload>/comm_dep.csv \
  --out out/<workload>_composed.csv \
  --l-min 0 --l-max 1000000 --step 50000
python3 pipeline/run_lgs.py \
  --goal <workload>/output.goal --L 1000
\end{verbatim}
\endgroup

The Composite-LP stage requires a Gurobi 10.0+ installation; the
LGS stage does not.

\emph{Precomputed outputs for resource-intensive stages.} Several
upstream stages of the pipeline are too expensive for a typical
reviewer machine --- in particular, LogGOPSim replay of the Grok
4{,}096-GPU trace, and the Monolithic-LP solves used as baselines
in Figures 5 and 7 (the largest of which requires several hundred
GiB of RAM). To make figure reproduction independent of access to
such machines, we ship the \emph{outputs} of these runs as
precomputed CSVs inside $A_1$: the LP sensitivity sweeps, the LGS
points, the hardware reference measurements, and the Grok memory
footprints all live under \texttt{data/} and are consumed directly
by the figure scripts. Reviewers can therefore validate all paper
figures from the packaged artifact without re-executing any LP
solve or packet-level simulation.

\artout

Running \texttt{reproduce\_all.py} writes the following files to
\texttt{figures/}:

\begin{itemize}
\item \texttt{fig\_2d\_sensitivity\_workloads.pdf} --- paper Figure~1.
\item \texttt{fig3\_sensitivity\_1x4.pdf} --- paper Figure~3.
\item \texttt{fig\_mixed\_16n\_ch1.pdf} --- paper Figure~4.
\item \texttt{fig5\_llama7b.pdf} --- paper Figure~5.
\item \texttt{fig\_3x3\_sensitivity.pdf} --- paper Figure~6.
\item \texttt{fig6\_grok\_memory.pdf} --- paper Figure~7.
\item \texttt{fig\_network\_perf\_combined.pdf} --- paper Figures 8 and~9 (combined).
\item \texttt{fig\_jitter\_3panel.pdf} --- paper Figure~10.
\end{itemize}

Each generated PDF should be visually identical to the corresponding
figure in the paper. A reviewer may perform the comparison by
opening each PDF next to the paper manuscript; the script
\texttt{reproduce\_all.py --list} prints the figure-to-script mapping
for easy cross-referencing.

\newartifact

\artrel

Artifact $A_2$ is the complete collection of NCCL execution traces
used throughout the paper, converted to the GOAL format, and
covering Llama 3.3 training (16--128 GPUs), Grok~3 training
(16--4{,}096 GPUs), and Llama 3.1-8B inference under vLLM (8 GPUs). By
publishing this dataset we fulfill contribution $C_1$: enabling
open, reproducible trace-driven network-sensitivity research.

\artexp

Each GOAL file in $A_2$ can be replayed with LogGOPSim or any other
GOAL-compatible simulator, and is the direct input to the
Composite-LP toolchain in $A_1$. The expected output of replaying a
trace is the corresponding $T(L)$ or $T(G)$ curve that appears in
the packaged CSVs of $A_1$. Table~II of the paper reports reference
latency and bandwidth values drawn from these traces.

\arttime

No reproduction time applies: $A_2$ is a data release. Download
time dominates and depends on the user's bandwidth (2\,TB total;
individual traces range from a few MiB for vLLM inference to
${\sim}1.3$\,TiB for Grok 314B at 4{,}096 GPUs).

\artin

\artinpart{Hardware}

None. $A_2$ is a data archive. Replaying its GOAL files end-to-end
under a packet-level simulator requires an x86 host with enough RAM
to hold the target trace (up to ${\sim}200$\,GiB for the largest
Grok run); replaying under LogGOPSim or our Composite-LP toolchain
requires only a few GiB.

\artinpart{Software}

A GOAL-compatible trace reader. All scripts in $A_1$ that need to
re-process traces use the reader shipped with the LogGOPSim sources
referenced in the main paper.

\artinpart{Datasets / Inputs}

Self-contained data release.

\artinpart{Installation and Deployment}

Only a subset of the traces published under
\texttt{http://storage2.spcl.ethz.ch/\allowbreak traces/ai/}
is consumed by the figures of this paper: Llama~3.3 training
at 16--128~GPUs, Grok~3 training at 16--4{,}096~GPUs, and Llama~3.1-8B
inference under vLLM at 8~GPUs (listed in Table~II of the
paper). We therefore recommend fetching only those specific
per-workload sub-trees rather than mirroring the full listing.
Example for a single workload:

\begingroup\small
\begin{verbatim}
wget -r -np -nH --cut-dirs=2 \
  http://storage2.spcl.ethz.ch/traces/ai/grok3/
\end{verbatim}
\endgroup

No build step is required.

\artcomp

$A_2$ is a static data release; there is no execution step. The
intended downstream workflow is \texttt{(trace $\in A_2$) $\to$ LP
solver (shipped in $A_1$) $\to$ sensitivity CSV (also shipped in
$A_1$) $\to$ figure}. Only the last two stages are exercised by the
reproduction workflow in $A_1$.

\artout

The artifact itself is the output. File integrity can be verified
against the SHA-256 manifest shipped alongside the trace archive.

\end{document}
